\documentclass[10pt,a4paper]{amsart}
\usepackage[margin=19mm]{geometry}
\usepackage{amsmath,amssymb,mathtools}
\usepackage[scr=rsfs,cal=euler]{mathalfa}
\usepackage{xfrac}
\usepackage[unicode,hidelinks,bookmarks=false]{hyperref}
\newcommand{\ie}{i.e.\ }
\newcommand{\cf}{cf.\ }
\newcommand{\resp}{respectively\ }
\newcommand{\Subsection}{\subsection}
\providecommand{\nobreakdash}{\nobreak\hbox{-}}
\setlength{\emergencystretch}{2em}
\multlinegap0pt
\usepackage[shortlabels]{enumitem}
\usepackage{amssymb}
\usepackage{mathtools}
\usepackage{bm}
\newtheorem{definition}{Definition}
\newtheorem{assumption}{Assumption}
\newtheorem{theorem}{Theorem}
\newtheorem{lemma}{Lemma}
\newtheorem{proposition}{Proposition}
\title{The circular law for non-Hermitian random band matrices up to bandwidth $N^{2/3+c}$}
\author{Yi Han}
\date{Source: arXiv:2508.18143v3 (CC BY 4.0)}
\begin{document}
\maketitle

\noindent\emph{Source and licence.} The circular law for non-Hermitian random band matrices up to bandwidth $N^{2/3+c}$. Version arXiv:2508.18143v3 (CC BY 4.0), distributed by arXiv under the Creative Commons Attribution 4.0 International licence, http://creativecommons.org/licenses/by/4.0/, as recorded in the arXiv licence metadata for this exact version and shown on the abstract page https://arxiv.org/abs/2508.18143v3. Full licence notice: \texttt{licenses/arxiv-2508.18143-CC-BY-4.0.txt}.\par\smallskip

\noindent\textit{Editorial scope.} Counted unit: the article's proposition prop:global (source0.tex lines 1347--1370). The article's complete original proof of it is delivered with it (source0.tex lines 1372--1533, one \texttt{proof} environment of 162 lines). No mathematics is written, repaired or re-derived: the statement and the proof are the article's own lines, in the article's own notation, and no retained line is substituted or re-flowed.  Retained uncounted as the source-local context the proof needs: lines 280--293; lines 295--305; lines 599--599; lines 638--661; lines 754--779; lines 960--972; lines 1094--1098; lines 1105--1107.  Nothing was omitted from any retained span, so the package carries no omission record.  No retained line contains a citation, so the wrapper carries no bibliography and no bibliography entry.  No retained line contains a figure environment, an image-inclusion command or drawing code, so no image is delivered and the package declares no asset dependency.  Editorial formatting only, in addition: an AMS \texttt{amsart} preamble that restates the article's own declarations and macros used by the retained lines, namely the environments \texttt{definition}, \texttt{assumption}, \texttt{theorem}, \texttt{lemma}, \texttt{proposition} on separate counters, together with the standard packages those lines need.  The retained environments print in this document as Definition 1, Assumption 1, Theorem 1, Lemma 1, Proposition 1 in order of appearance, whereas the article prints the same items with the numbering of its own full document.  The complete native source of the article as distributed in the arXiv e-print for this exact version is bundled unchanged as \texttt{sources/183837/source0.tex}.
\begin{definition}[Normalized independent-entry model]\label{mainassumptionbands}
Let $S=(s_{ij})_{i,j=1}^N$ be a deterministic nonnegative matrix. Set
\[
 X_{ij}=b_{ij}\xi_{ij},\qquad b_{ij}=\sqrt{s_{ij}},
 \qquad \mathbb E\xi_{ij}=0,\quad \mathbb E|\xi_{ij}|^2=1,
\]
where the variables $\xi_{ij}$ are independent. We assume
\begin{equation}\label{eq:normalized-profile}
 \sum_j s_{ij}=\sum_j s_{ji}=1,\qquad
 s_*:=\max_{i,j}s_{ij}\leq \frac{C_0}{W}.
\end{equation}
Here $W=W_N\geq1$ and $C_0$ is independent of $N$. Thus
$W\leq C_0N$. Zero variances are allowed. 
\end{definition}

\begin{assumption}[A profile symmetry ensuring balance]\label{ass:profile-balance}
The variance profile in Definition~\ref{mainassumptionbands} satisfies at
least one of the following conditions:
\begin{enumerate}[label=\textup{(\roman*)}]
\item $S=S^T$;
\item there is a transitive subgroup $\mathfrak G_N$ of the permutations
of $[N]$ such that $s_{\pi(i),\pi(j)}=s_{ij}$ for all $\pi\in\mathfrak G_N$.
\end{enumerate}
Transitive means that any index can be sent to any other index by a member
of $\mathfrak G_N$; the group is allowed to depend on $N$.
\end{assumption}

\section{A coarse Gaussian resolvent estimate}\label{sec:coarse}

\begin{theorem}[Coarse expected resolvent bound]\label{thm:coarse}
There are universal constants $c,C>0$ such that, under the assumptions of
this section, for every $z\ne0$ and every $\eta>0$ satisfying
\begin{equation}\label{eq:coarse-scale-condition}
 \frac{s_*}{\eta^2}\le c,
\end{equation}
one has
\begin{equation}\label{eq:coarse-coordinate-bound}
 \mathbb E a_i^+=\mathbb E a_i^-\le\frac{C}{|z|}
 \qquad (1\le i\le N).
\end{equation}
In particular,
\begin{equation}\label{eq:coarse-trace-bound}
 \mathbb E\frac1N\sum_{\alpha=1}^N
 \frac{\eta}{s_\alpha(X-zI_N)^2+\eta^2}
 \le\frac{C}{|z|}.
\end{equation}
If $s_*\le C_0/W$, these conclusions hold for every
$\eta\ge K W^{-1/2}$, where $K$ depends only on $C_0$.
More generally, at a fixed $z,\eta$ the same conclusions hold for any
Gaussian variance profile for which
$\mathbb E a_i^+=\mathbb E a_i^-$ for every $i$; the two structural
profile conditions above are sufficient ways to establish this balance.
\end{theorem}

\begin{lemma}\label{lem:coarse-ward}
The block entries of $R$ satisfy
\begin{align}
 R_{++}&=i\eta(YY^*+\eta^2I_N)^{-1},&
 R_{--}&=i\eta(Y^*Y+\eta^2I_N)^{-1},
 \label{eq:coarse-diagonal-blocks}\\
 R_{+-}&=Y(Y^*Y+\eta^2I_N)^{-1},&
 R_{-+}&=Y^*(YY^*+\eta^2I_N)^{-1}=R_{+-}^*.
 \label{eq:coarse-offdiagonal-blocks}
\end{align}
In particular, $0\le a_i^\pm\le\eta^{-1}$ and
$R_{i-,i+}=\overline{t_i}$.  For every Hermitization index $v$,
\begin{equation}\label{eq:coarse-ward}
 \sum_u|R_{vu}|^2=\sum_u|R_{uv}|^2
 =\frac{\operatorname{Im}R_{vv}}{\eta}.
\end{equation}
Moreover,
\begin{equation}\label{eq:coarse-transpose-means}
 \mathbb E a_i^+=\mathbb E a_i^-=:A_i.
\end{equation}
For every variance profile, independently of either structural condition,
one has the pathwise identity
\begin{equation}\label{eq:coarse-exact-trace-balance}
 \sum_i a_i^+=\sum_i a_i^-.
\end{equation}
\end{lemma}

\begin{lemma}\label{lem:coarse-expectation-equations}
If $\rho\le1/2$, then
\begin{align}
 1&=A_iB_i-z\overline{\tau_i}+e_{1i},
 \label{eq:coarse-SD-diagonal}\\
 0&=B_i\tau_i+zA_i+e_{2i},
 \label{eq:coarse-SD-offdiagonal}
\end{align}
where
\begin{equation}\label{eq:coarse-error-bound}
 |e_{1i}|+|e_{2i}|\le C\rho A_iB_i.
\end{equation}
\end{lemma}

\begin{equation}\label{eq:coarse-scalar-equation}
 1=A_iB_i+\frac{|z|^2A_i}{B_i}+\mathcal E_i,
 \qquad
 \mathcal E_i=e_{1i}+\frac{z\overline{e_{2i}}}{B_i}.
\end{equation}

\begin{equation}\label{eq:coarse-scalar-error}
 |\mathcal E_i|\le C\rho A_iB_i+C\rho|z|A_i.
\end{equation}

\begin{proposition}\label{prop:global}
Let $X_N$ have independent centered real Gaussian entries, or independent
centered circular complex Gaussian entries, with a doubly stochastic
variance profile $S_N$ satisfying Assumption~\ref{ass:profile-balance}
(symmetric or invariant under a transitive group of simultaneous index
permutations).  Suppose
$s_{*,N}:=\max_{i,j}(S_N)_{ij}\to0$.  For every fixed $z\in\mathbb C$, the
empirical singular-value measure
\[
 \nu_{N,z}:=\frac1N\sum_{\alpha=1}^N
 \delta_{s_\alpha(X_N-zI_N)}
\]
converges weakly in probability to a deterministic probability measure
$\nu_z$ on $[0,\infty)$.  This measure is the same for both Gaussian
classes and for every such sequence of variance profiles.  In particular,
if $G_N$ is a normalized dense Ginibre matrix of either class, with
$\mathbb E|(G_N)_{ij}|^2=N^{-1}$, then for every bounded continuous
$f:[0,\infty)\to\mathbb R$,
\begin{equation}\label{eq:global-comparison}
 \int f\,d\nu_{N,z}-\int f\,d\nu_{G_N-zI_N}
 \longrightarrow0\qquad\text{in probability}.
\end{equation}
The matrices in this comparison need not be independent.
\end{proposition}

\begin{proof}
We retain the notation of Section~\ref{sec:coarse};
Lemma~\ref{lem:coarse-ward} gives the diagonal expectation balance under
either profile condition.  First fix
$\eta\in(2,3)$.  The deterministic resolvent bound and stochasticity give
\begin{equation}\label{eq:global-A-B-bounds}
 0\le A_i\le\eta^{-1},\qquad
 \eta\le B_i\le\eta+\eta^{-1}.
\end{equation}
For all sufficiently large $N$, the assumption
$\rho=4s_{*,N}/\eta^2\le1/2$ holds.  Thus the exact equations and error
bounds in Lemma~\ref{lem:coarse-expectation-equations} apply.  In
particular, \eqref{eq:coarse-scalar-equation}--\eqref{eq:coarse-scalar-error}
and \eqref{eq:global-A-B-bounds} imply
\begin{equation}\label{eq:global-approx-fixed-point}
 A_i=f_\eta((S_NA)_i)+O_z(s_{*,N}),\qquad
 f_\eta(x):=\frac{\eta+x}{(\eta+x)^2+|z|^2}.
\end{equation}
The error is uniform in $i$ and in $\eta\in(2,3)$.  Indeed,
$A_i=f_\eta((S_NA)_i)(1-\mathcal E_i)$ by
\eqref{eq:coarse-scalar-equation}, and
$|\mathcal E_i|\le C_z\,s_{*,N}$ on this fixed-scale interval.  This
argument also applies at $z=0$: it does not use the bound $C/|z|$ in
Theorem~\ref{thm:coarse}.

The map $f_\eta$ sends $[0,\eta^{-1}]$ into itself and satisfies
\[
 |f_\eta'(x)|
 =\frac{\bigl||z|^2-(\eta+x)^2\bigr|}
        {((\eta+x)^2+|z|^2)^2}
 \le\frac1{(\eta+x)^2}\le\eta^{-2}<1.
\]
Consequently there is a unique fixed point
$a_z(\eta)\in[0,\eta^{-1}]$ such that
\begin{equation}\label{eq:global-scalar-fixed-point}
 a_z(\eta)=\frac{\eta+a_z(\eta)}
 {(\eta+a_z(\eta))^2+|z|^2}.
\end{equation}
Since $S_N\mathbf1=\mathbf1$ and
$\|S_N\|_{\ell^\infty\to\ell^\infty}=1$, subtracting the constant fixed
point from \eqref{eq:global-approx-fixed-point} gives
\[
 \max_i|A_i-a_z(\eta)|
 \le\eta^{-2}\max_i|A_i-a_z(\eta)|+C_z\,s_{*,N}.
\]
It follows that
\begin{equation}\label{eq:global-mean-convergence}
 \max_i|A_i-a_z(\eta)|\le C_z\,s_{*,N}.
\end{equation}
This is a large-$\eta$ contraction only, not a stability statement down to
a vanishing imaginary part.

We also need concentration of the normalized trace.  Set
\[
 m_N(i\eta):=\frac1{2N}\operatorname{Tr}R_z(\eta)
 =i\int\frac{\eta}{s^2+\eta^2}\,d\nu_{N,z}(s).
\]
The equality follows because $YY^*$ and $Y^*Y$ have the same eigenvalues.
In the circular complex case,
\[
 \partial_{x_{ab}}m_N(i\eta)
 =-\frac{(R^2)_{b-,a+}}{2N},\qquad
 \partial_{\overline{x}_{ab}}m_N(i\eta)
 =-\frac{(R^2)_{a+,b-}}{2N}.
\]
In the real case the derivative is the sum of these two expressions.
Gaussian Poincar\'e therefore gives in both cases
\begin{equation}\label{eq:global-trace-variance}
 \operatorname{Var}(m_N(i\eta))
 \le\frac{C_{\mathrm{Poincar\acute e}}\,s_{*,N}}{N^2}\mathbb E\|R^2\|_{\mathrm{HS}}^2
 \le\frac{C_{\mathrm{Poincar\acute e}}\,s_{*,N}}{N\eta^4}.
\end{equation}
Here $C_{\mathrm{Poincar\acute e}}$ is an absolute constant, independent of $N$, $z$, $\eta$,
and the variance profile, and we recall that
$s_{*,N}=\max_{a,b}(S_N)_{ab}$. Thus $C_{\mathrm{Poincar\acute e}}\,s_{*,N}$ denotes the
product of this constant and the maximum entry variance.
The factor $s_{*,N}$ comes from bounding each variance weight in the
Gaussian Poincar\'e inequality by $s_{*,N}$.
The last inequality uses the pathwise bound
$\|R^2\|_{\mathrm{HS}}^2\le2N\eta^{-4}$.  Combining
\eqref{eq:global-mean-convergence} and
\eqref{eq:global-trace-variance} yields, for each $\eta\in(2,3)$,
\begin{equation}\label{eq:global-resolvent-test}
 \int h_\eta(s)\,d\nu_{N,z}(s)
 \longrightarrow a_z(\eta)\quad\text{in probability},
 \qquad h_\eta(s):=\frac{\eta}{s^2+\eta^2}.
\end{equation}
The same convergence holds for expectations.

We spell out why these resolvent tests determine weak convergence.
First, centering and double stochasticity imply
\begin{equation}\label{eq:global-second-moment}
 \mathbb E\int s^2\,d\nu_{N,z}(s)
 =\frac1N\mathbb E\|X_N-zI_N\|_{\mathrm{HS}}^2
 =1+|z|^2.
\end{equation}
Thus the deterministic probability measures
$\overline\nu_{N,z}:=\mathbb E\nu_{N,z}$ are tight.  Any subsequential
weak limit $\nu$ satisfies
\begin{equation}\label{eq:global-measure-characterization}
 \int h_\eta\,d\nu=a_z(\eta)
 \qquad\bigl(\eta\in\mathbb Q\cap(2,3)\bigr),
\end{equation}
because $h_\eta$ is bounded and continuous.

There is at most one probability measure on $[0,\infty)$ satisfying
\eqref{eq:global-measure-characterization}. To prove this, let $\nu_1$
and $\nu_2$ be two such measures. Define $T:[0,\infty)\to[0,\infty)$
by $T(s)=s^2$, and let $\widetilde\nu_j:=T_\#\nu_j$ be the
pushforward measure, that is,
\[
 \widetilde\nu_j(B)=\nu_j\bigl(T^{-1}(B)\bigr)
 \quad\text{for every Borel set }B\subset[0,\infty),\qquad j=1,2.
\]
Their Stieltjes transforms are
\[
 q_j(w):=\int\frac1{t-w}\,d\widetilde\nu_j(t)
       =\int\frac1{s^2-w}\,d\nu_j(s),
 \qquad w\in\mathbb C\setminus[0,\infty).
\]
For every rational $\eta\in(2,3)$,
\[
 q_j(-\eta^2)=\int\frac1{s^2+\eta^2}\,d\nu_j(s)
 =\frac1\eta\int h_\eta(s)\,d\nu_j(s)
 =\frac{a_z(\eta)}\eta,
\]
so $q_1$ and $q_2$ agree at all such points $-\eta^2$. These points
have an accumulation point inside the connected domain
$\mathbb C\setminus[0,\infty)$. The identity theorem gives $q_1=q_2$
throughout this domain, and uniqueness of the Stieltjes transform gives
$\widetilde\nu_1=\widetilde\nu_2$. Since $T$ is a homeomorphism of
$[0,\infty)$ onto itself, it follows that $\nu_1=\nu_2$.
Tightness supplies existence, so denote the resulting unique measure by
$\nu_z$.

Finally, regard $\nu_{N,z}$ as a random element of the space of probability
measures on $[0,\infty)$ with its weak topology.  For $M>0$, the set
\[
 \mathcal K_M:=\left\{\nu:\int s^2\,d\nu(s)\le M\right\}
\]
is compact: its measures are uniformly tight, and it is closed by the
lower semicontinuity of the second moment.  Markov's inequality and
\eqref{eq:global-second-moment} give
\[
 \mathbb P\{\nu_{N,z}\notin\mathcal K_M\}
 \le\frac{1+|z|^2}{M}.
\]
Hence the laws of these random measures are tight.  Every subsequential
limit in distribution satisfies
\eqref{eq:global-measure-characterization} almost surely, by
\eqref{eq:global-resolvent-test}, the continuity of integration against
each $h_\eta$, and the countability of $\mathbb Q\cap(2,3)$.  By the
uniqueness just proved, that limit equals $\nu_z$ almost surely.
Convergence in distribution to this deterministic measure is convergence
in probability.  This proves the first claim.

The normalized dense real and circular complex Ginibre matrices satisfy
the same assumptions, with $s_{*,N}=N^{-1}$.  The scalar fixed point and
the characterizing tests are identical.  Applying the result to either
Ginibre class and taking the difference of the two convergences proves
\eqref{eq:global-comparison}, regardless of the coupling of the matrices.
\end{proof}

\end{document}
